\documentclass[11pt]{article}
\newcommand{\ListaTitulo}{Gabarito 06: Probabilidade condicional e independência}
% Preâmbulo compartilhado: MATF14, Listas de Exercícios 2026
% Usado por ListaNN.tex e GabaritoNN.tex via % Preâmbulo compartilhado: MATF14, Listas de Exercícios 2026
% Usado por ListaNN.tex e GabaritoNN.tex via % Preâmbulo compartilhado: MATF14, Listas de Exercícios 2026
% Usado por ListaNN.tex e GabaritoNN.tex via \input{preamble}

\usepackage[utf8]{inputenc}
\usepackage[T1]{fontenc}
\usepackage[brazil]{babel}
\usepackage{fullpage}
\usepackage{amsmath,amssymb}
\usepackage{enumitem}
\usepackage{xcolor}
\usepackage{fancyhdr}
\usepackage{titlesec}
\usepackage{listings}
\usepackage{hyperref}
\usepackage{booktabs}
\usepackage{graphicx}
\usepackage{tikz}

\definecolor{matf14azul}{HTML}{0D3B54}
\definecolor{matf14dourado}{HTML}{C9971E}

\titleformat{\section}{\normalfont\Large\bfseries\color{matf14azul}}{\thesection}{1em}{}
\titleformat{\subsection}{\normalfont\large\bfseries\color{matf14azul}}{\thesubsection}{1em}{}

\providecommand{\ListaTitulo}{Lista de Exercícios}

\setlength{\headheight}{14pt}
\pagestyle{fancy}
\fancyhf{}
\lhead{\small MATF14: Estatística Econômica I}
\rhead{\small \ListaTitulo}
\cfoot{\thepage}
\renewcommand{\headrulewidth}{0.4pt}

\lstset{
  basicstyle=\ttfamily\small,
  breaklines=true,
  frame=single,
  columns=fullflexible,
  backgroundcolor=\color{gray!8},
  commentstyle=\color{matf14dourado},
  keywordstyle=\color{matf14azul}\bfseries,
  language=R
}

\newcounter{questaocounter}
\newenvironment{questao}[1]{%
  \stepcounter{questaocounter}%
  \bigskip\noindent\textbf{Questão #1.}\ %
}{\par}

\newcommand{\cabecalho}[2]{%
  \begin{center}
    {\Large\bfseries\color{matf14azul} #1}\\[0.3em]
    {\large #2}
  \end{center}
  \vspace{0.5em}
  \noindent\rule{\linewidth}{0.6pt}
  \vspace{0.5em}
}

\newenvironment{solucao}{%
  \par\smallskip\noindent\textit{\color{matf14dourado!80!black}Solução.}\ %
}{\hfill$\blacksquare$\par}


\usepackage[utf8]{inputenc}
\usepackage[T1]{fontenc}
\usepackage[brazil]{babel}
\usepackage{fullpage}
\usepackage{amsmath,amssymb}
\usepackage{enumitem}
\usepackage{xcolor}
\usepackage{fancyhdr}
\usepackage{titlesec}
\usepackage{listings}
\usepackage{hyperref}
\usepackage{booktabs}
\usepackage{graphicx}
\usepackage{tikz}

\definecolor{matf14azul}{HTML}{0D3B54}
\definecolor{matf14dourado}{HTML}{C9971E}

\titleformat{\section}{\normalfont\Large\bfseries\color{matf14azul}}{\thesection}{1em}{}
\titleformat{\subsection}{\normalfont\large\bfseries\color{matf14azul}}{\thesubsection}{1em}{}

\providecommand{\ListaTitulo}{Lista de Exercícios}

\setlength{\headheight}{14pt}
\pagestyle{fancy}
\fancyhf{}
\lhead{\small MATF14: Estatística Econômica I}
\rhead{\small \ListaTitulo}
\cfoot{\thepage}
\renewcommand{\headrulewidth}{0.4pt}

\lstset{
  basicstyle=\ttfamily\small,
  breaklines=true,
  frame=single,
  columns=fullflexible,
  backgroundcolor=\color{gray!8},
  commentstyle=\color{matf14dourado},
  keywordstyle=\color{matf14azul}\bfseries,
  language=R
}

\newcounter{questaocounter}
\newenvironment{questao}[1]{%
  \stepcounter{questaocounter}%
  \bigskip\noindent\textbf{Questão #1.}\ %
}{\par}

\newcommand{\cabecalho}[2]{%
  \begin{center}
    {\Large\bfseries\color{matf14azul} #1}\\[0.3em]
    {\large #2}
  \end{center}
  \vspace{0.5em}
  \noindent\rule{\linewidth}{0.6pt}
  \vspace{0.5em}
}

\newenvironment{solucao}{%
  \par\smallskip\noindent\textit{\color{matf14dourado!80!black}Solução.}\ %
}{\hfill$\blacksquare$\par}


\usepackage[utf8]{inputenc}
\usepackage[T1]{fontenc}
\usepackage[brazil]{babel}
\usepackage{fullpage}
\usepackage{amsmath,amssymb}
\usepackage{enumitem}
\usepackage{xcolor}
\usepackage{fancyhdr}
\usepackage{titlesec}
\usepackage{listings}
\usepackage{hyperref}
\usepackage{booktabs}
\usepackage{graphicx}
\usepackage{tikz}

\definecolor{matf14azul}{HTML}{0D3B54}
\definecolor{matf14dourado}{HTML}{C9971E}

\titleformat{\section}{\normalfont\Large\bfseries\color{matf14azul}}{\thesection}{1em}{}
\titleformat{\subsection}{\normalfont\large\bfseries\color{matf14azul}}{\thesubsection}{1em}{}

\providecommand{\ListaTitulo}{Lista de Exercícios}

\setlength{\headheight}{14pt}
\pagestyle{fancy}
\fancyhf{}
\lhead{\small MATF14: Estatística Econômica I}
\rhead{\small \ListaTitulo}
\cfoot{\thepage}
\renewcommand{\headrulewidth}{0.4pt}

\lstset{
  basicstyle=\ttfamily\small,
  breaklines=true,
  frame=single,
  columns=fullflexible,
  backgroundcolor=\color{gray!8},
  commentstyle=\color{matf14dourado},
  keywordstyle=\color{matf14azul}\bfseries,
  language=R
}

\newcounter{questaocounter}
\newenvironment{questao}[1]{%
  \stepcounter{questaocounter}%
  \bigskip\noindent\textbf{Questão #1.}\ %
}{\par}

\newcommand{\cabecalho}[2]{%
  \begin{center}
    {\Large\bfseries\color{matf14azul} #1}\\[0.3em]
    {\large #2}
  \end{center}
  \vspace{0.5em}
  \noindent\rule{\linewidth}{0.6pt}
  \vspace{0.5em}
}

\newenvironment{solucao}{%
  \par\smallskip\noindent\textit{\color{matf14dourado!80!black}Solução.}\ %
}{\hfill$\blacksquare$\par}


\begin{document}

\cabecalho{Gabarito 06}{Probabilidade condicional e independência estatística (Cap. 2, Aulas 11--12)}

\begin{questao}{1} \textbf{(Condicional: inadimplência por faixa de renda)}
\begin{solucao}
\begin{enumerate}[label=(\alph*)]
  \item $P(\text{inad}\mid\text{baixa}) = 80/400 = 0{,}20$.
  \item $P(\text{inad}\mid\text{alta}) = 5/200 = 0{,}025$.
  \item $P(\text{inad})$ geral $=105/1000=0{,}105$. As três probabilidades condicionais
    ($0{,}20$ para renda baixa, $5{,}0\%$ para renda média, $20/400$, e $2{,}5\%$ para renda alta)
    são bem diferentes entre si e diferentes da probabilidade geral, \textbf{não há} independência:
    a faixa de renda carrega informação relevante sobre a probabilidade de inadimplência nesta
    base.
\end{enumerate}
\end{solucao}
\end{questao}

\begin{questao}{2} \textbf{(Regra do produto: sorteio sem reposição)}
\begin{solucao}
\begin{enumerate}[label=(\alph*)]
  \item $P(1^\circ\text{ irregular}) = 4/15$. Dado que o primeiro foi irregular, sobram 3
    irregulares em 14: $P(2^\circ\text{ irregular}\mid 1^\circ\text{ irregular}) = 3/14$. Logo,
    $P(\text{ambos irregulares}) = \frac{4}{15}\times\frac{3}{14} = \frac{12}{210} = \frac{2}{35}
    \approx 0{,}057$.
  \item $P(1^\circ\text{ regular})=11/15$; dado isso, sobram 10 regulares em 14:
    $P(2^\circ\text{ regular}\mid 1^\circ\text{ regular})=10/14$. Logo,
    $P(\text{nenhum irregular}) = \frac{11}{15}\times\frac{10}{14} = \frac{110}{210} =
    \frac{11}{21} \approx 0{,}524$.
  \item Com reposição: $P(\text{ambos irregulares}) = \frac{4}{15}\times\frac{4}{15} =
    \frac{16}{225}\approx 0{,}071$, maior que o valor sem reposição ($0{,}057$). Sem reposição, o
    primeiro sorteio "consome" um processo irregular, reduzindo a proporção de irregulares
    disponíveis para o segundo sorteio ($3/14 < 4/15$), com reposição, a proporção permanece
    constante em $4/15$ nos dois sorteios, sendo eventos independentes.
\end{enumerate}
\end{solucao}
\end{questao}

\begin{questao}{3} \textbf{(Independência: verificação numérica)}
\begin{solucao}
\begin{enumerate}[label=(\alph*)]
  \item $P(\text{comércio}) = 120/200 = 0{,}6$. $P(\text{crédito}) = 60/200 = 0{,}3$.
    $P(\text{comércio}\cap\text{crédito}) = 36/200 = 0{,}18$.
  \item $P(\text{comércio})\times P(\text{crédito}) = 0{,}6\times 0{,}3 = 0{,}18$, exatamente igual
    a $P(\text{comércio}\cap\text{crédito})=0{,}18$ observado. Nesta amostra, setor (comércio ou
    não) e acesso a crédito subsidiado \textbf{são independentes}, a proporção de empresas de
    comércio é a mesma dentro e fora do grupo que acessou crédito.
\end{enumerate}
\end{solucao}
\end{questao}

\begin{questao}{4} \textbf{(Mutuamente exclusivos vs. independentes: discussão numérica)}
\begin{solucao}
\begin{enumerate}[label=(\alph*)]
  \item Mutuamente exclusivos $\Rightarrow P(A\cap B)=0$ (por definição). Logo,
    $P(A\mid B) = P(A\cap B)/P(B) = 0/0{,}5 = 0$.
  \item Independentes $\Rightarrow P(A\cap B) = P(A)P(B) = 0{,}4\times 0{,}5 = 0{,}20$. Logo,
    $P(A\mid B) = P(A) = 0{,}4$ (por definição de independência).
  \item Não é possível. Se fossem mutuamente exclusivos, $P(A\cap B)$ teria que ser $0$; se fossem
    independentes (com essas probabilidades individuais positivas), $P(A\cap B)$ teria que ser
    $0{,}20 \ne 0$. Como $P(A)>0$ e $P(B)>0$, os dois valores exigidos são incompatíveis entre si, dois eventos de probabilidade positiva não podem ser simultaneamente mutuamente exclusivos e
    independentes.
\end{enumerate}
\end{solucao}
\end{questao}

\begin{questao}{5} \textbf{(Dedução: independência é simétrica)}
\begin{solucao}
\begin{enumerate}[label=(\alph*)]
  \item Pela definição de probabilidade condicional, $P(A\mid B) = \dfrac{P(A\cap B)}{P(B)}$.
    Substituindo a hipótese $P(A\cap B)=P(A)P(B)$:
    $$
    P(A\mid B) = \frac{P(A)P(B)}{P(B)} = P(A),
    $$
    usando que $P(B)>0$ para poder cancelar.
  \item Analogamente, $P(B\mid A) = \dfrac{P(A\cap B)}{P(A)} = \dfrac{P(A)P(B)}{P(A)} = P(B)$,
    usando $P(A)>0$. Como a hipótese $P(A\cap B)=P(A)P(B)$ é simétrica em $A$ e $B$ (o produto não
    depende da ordem), as duas conclusões, $P(A\mid B)=P(A)$ e $P(B\mid A)=P(B)$, decorrem da
    mesma hipótese única, confirmando que a independência de $A$ em relação a $B$ é equivalente à
    independência de $B$ em relação a $A$.
\end{enumerate}
\end{solucao}
\end{questao}

\end{document}
